\clearpage
\TemplateRow[1.5pt][.75pt]{结\quad 论\\（结\quad 果）}{
仿真网络完成A/CNAME配置与三阶段测试。初始三种ping均4/4；DNS服务器地址改变而PC仍用旧地址时，IP访问保持成功、两个域名均失败；修正PC DNS后全部恢复4/4。A记录始终指向目标主机192.168.1.200。别名响应ANCOUNT=2，但该版本界面只展开末端A块，已据实限定字段证据范围。\par 
本轮真实抓包完成两个公网域名首次/重复ping与四种nslookup，最终文件的6组DNS事务按ID、问题、反向端点和时间共同配对。百度存在指向\domain{www.a.shifen.com}的CNAME；川大本次返回直接A，CNAME专门查询为NOERROR/NODATA。
\par 系统DNS与显式公网解析器对百度给出不同IPv4地址，分别记录了各自结果和实际ping目标。首次有DNS、重复无DNS且保持相同地址的现象支持缓存复用。川大虽然ICMP未响应，但成功A回答已证明名称解析完成。
}{94mm}
\TemplateRow{小\quad 结}{
本次实验使我能够把IP连通性、名称解析和应用测试分开判断。查看完整Answer及各分区后，才能区分CNAME、网站A地址和SOA名称服务器；根据QR与反向端点等共同配对，也避免了只凭事务ID关联报文。
\par macOS缓存条目转储未能成功，本次仅通过清缓存命令状态、首次与重复报文变化分析缓存；代理和远端缓存未清理。抓包范围也不包含远端解析器完整上游过程，不能扩展为完整递归链路证据。
\par 对川大ICMP无响应，仅能记录现象，不能直接归因于DNS错误或断定服务器禁ping。本轮没有调整系统网络设置，无需恢复临时DNS配置。后续复测应保持相同解析器、查询类型与采集窗口，并继续以原始报文核对命令输出。
}{83mm}
\TemplateReview
